\emph{(i).} Under the flow, $\alpha_{t}(a^{*}) = e^{i\omega t}a^{*}$ and $\alpha_{t}(a) = e^{-i\omega t}a$, since $[K, a] = -a$ (Proposition \claimref{hawking.modular-structure-completed-pair}(2), (4)). With occupancy $n = \omega(a^{*}a)$ and $\omega(aa^{*}) = 1 - n$, the ladder correlators are $F_{a,a^{*}}(t) = e^{i\omega t}(1 - n)$ and $F_{a^{*},a}(t) = e^{-i\omega t}n$, and $F_{a,a} = F_{a^{*},a^{*}} = 0$ since $a^{2} = 0$. Each is entire, so regularity alone imposes nothing. Strict periodicity gives $(Q^{-1} - 1)(1 - n) = 0$ and $(Q - 1)n = 0$; for $Q \ne 1$ these force $n = 1$ and $n = 0$, so no state satisfies them. The glued reading gives $Q^{-1}(1 - n) = n$ (and the same equation from the other ordering), whose one root is $n = 1/(Q + 1)$, the occupancy T05's (R) asserts on the mode. The ladder correlators involve only $\omega(aa^{*})$ and $\omega(a^{*}a)$, so they leave $\omega(a)$ free; with $\omega(a) = 0$, which (P) (\claimref{singularity.p-univalence}) supplies, the marginal is $\operatorname{diag}(1 - n, n) = e^{-pK}/\operatorname{tr}e^{-pK}$, the unique KMS state at $\beta = p$ (Theorem \claimref{singularity.kms-uniqueness-factor}). At $Q = 1$ both periodicity equations hold for every $n$, and the glued one reads $1 - n = n$. The check runs these equations for $Q = 3, \tfrac15, \tfrac72$ with exact density matrices, and at $Q = 1$.

\emph{(ii).} For $i \ne j$ the canonical anticommutation relations (Proposition \claimref{scrambling.boundary-mode-space-matrix}) give $b_{j}a_{i} = -a_{i}b_{j}$ for $b_{j} \in \{a_{j}, a_{j}^{*}\}$. The glued condition for $(a_{i}, a_{j}^{*})$ reads $Q_{j}^{-1}\omega(a_{i}a_{j}^{*}) = \omega(a_{j}^{*}a_{i}) = -\omega(a_{i}a_{j}^{*})$, so $(1 + Q_{j}^{-1})\,\omega(a_{i}a_{j}^{*}) = 0$ and $\omega(a_{i}a_{j}^{*}) = 0$ since $Q_{j} > 0$; likewise for $(a_{i}, a_{j})$ with $Q_{j}$ in place of $Q_{j}^{-1}$, and for the adjoint pairs. Conversely, if every inter-mode two-point function vanishes, the mixed conditions hold trivially, and the single-mode ones are (i).

\emph{(iii).} The two-point clause of (Q) is $\omega(a_{k}^{*}a_{l}) = \delta_{kl}n_{k}$ with gauge invariance, so every inter-mode two-point function vanishes and, with T05's (R), (ii) gives (R$_{\mathrm{reg}}$); (R$_{\mathrm{reg}}$) implies T05's (R) by restriction to single modes. For the counter-models take $Q_{1} = 3$, $Q_{2} = 5$, $n_{k} = 1/(Q_{k} + 1)$, in the Jordan-Wigner basis. Model A is the product state with the coherence $c = \tfrac{1}{50}$ added between $\ket{01}$ and $\ket{10}$: it has trace one, is positive (its $2\times2$ block has determinant $n_{1}(1 - n_{2})\,n_{2}(1 - n_{1}) - c^{2} > 0$), commutes with the total parity, has the Gibbs marginals, and $\omega(a_{1}a_{2}^{*}) = -c$, so the mixed glued condition fails. Model B shifts $p_{11}$ to $n_{1}n_{2} + d$, $d = \tfrac{1}{100}$, with the other diagonal entries fixed by the marginals and no coherence: it is a diagonal, hence even, state with the Gibbs marginals and vanishing inter-mode two-point functions, so (R$_{\mathrm{reg}}$) holds, while $\omega(N_{1}N_{2}) - n_{1}n_{2} = d$, against Wick's factorisation $\omega(N_{1}N_{2}) = n_{1}n_{2}$ for a gauge-invariant quasi-free state with diagonal two-point function. The check confirms each statement exactly.

\emph{(iv).} By (i), (R$_{\mathrm{reg}}$) with (P) makes each marginal the Gibbs state at $\beta = p = 2\pi/\kappa$, with $n = 1/(e^{\omega p} + 1) \in (0, 1)$ for every real $\omega$, so every marginal is faithful; this is the single-mode conclusion of Theorem \claimref{singularity.rate-discharge-hypothesis-iii}, there drawn from T05's (R), and no modular group, hence no geometric modular condition, enters. The single-mode conclusions of the note's Theorem \claimref{state-selection.rate-discharge-hypothesis-iii} are these.

\emph{(v).} Both premises take the generators to be the boost bivector field and its products, so the generated algebra is the one of $\sigma_{1}, \ldots, \sigma_{n}$. The check builds three modes in $\Cl(3,3)$, one hyperbolic atom per mode: the $8$ products of subsets of $\{\sigma_{1}, \sigma_{2}, \sigma_{3}\}$ are linearly independent, each commutes with every $N_{k}$ (so every generator is diagonal in the joint number basis, as Proposition \claimref{singularity.number-diagonal-dynamics-models-scrambling-derived} states), $\sigma_{1}\sigma_{2}$ lies in their span, and the hopping $a_{1}^{*}a_{2}$ and pairing $a_{1}^{*}a_{2}^{*}$ each raise the rank to $9$ and fail to commute with $N_{1}$. T05's premise states the exclusion as no off-diagonal generator, in particular no hopping, with cross-point products allowed; the note's states it as no generator coupling two ladders. Since the note's own Proposition uses the cross-point controlled-$Z$, a product of $\sigma_{j}\sigma_{k}$-type terms, the literal reading that excludes all two-point generators contradicts its first clause, and the reading consistent with it excludes exactly the off-diagonal generators, which is T05's.
